\chapter{Introducción}
\label{cap:introduccion}

Presente el contexto general de la investigación, el problema y la estructura
del documento. Puede citar una fuente de ejemplo con \textcite{ejemplo2026}.

\section{Planteamiento del problema}

Describa con precisión el problema que motiva la investigación.

\begin{bloque}[Objetivo de este bloque]
  El entorno \texttt{bloque} permite destacar notas, resultados o ideas
  centrales sin mezclarlas con el texto principal.
\end{bloque}

\begin{definicion}[Concepto central]
  Use \texttt{definicion} para introducir términos que se emplearán a lo largo
  del trabajo.
\end{definicion}

\begin{advertencia}
  Reemplace todos los textos de ejemplo antes de entregar la tesis.
\end{advertencia}

\section{Objetivos}

\subsection{Objetivo general}

Formule el propósito principal de la investigación.

\subsection{Objetivos específicos}

\begin{itemize}
  \item Primer objetivo específico.
  \item Segundo objetivo específico.
  \item Tercer objetivo específico.
\end{itemize}
